\section{Компьютерная графика. Работа с графическими библиотеками}

\textbf{Компьютерная графика} изучает методы математического описания, создания, обработки и визуализации изображений средствами вычислительной техники. В задачах моделирования она используется прежде всего для представления геометрии и результатов расчёта.

Общая схема растеризационной визуализации имеет вид
\[
\text{модель сцены}
\rightarrow
\text{геометрические преобразования}
\rightarrow
\text{проекция}
\rightarrow
\text{растеризация}
\rightarrow
\text{изображение}.
\]

\subsection{Растровое и векторное представление}

\textbf{Растровое изображение} задаётся прямоугольной матрицей пикселей
\[
I=\{c_{ij}\},
\]
где $c_{ij}$ хранит цвет; например, в модели RGB
\[
c_{ij}=(R_{ij},G_{ij},B_{ij}),
\]
а дополнительный канал $\alpha$ может задавать прозрачность. Растровое изображение зависит от разрешения.

\textbf{Векторная графика} описывает изображение геометрическими примитивами: точками, отрезками, кривыми, многоугольниками, текстом. Она удобна для масштабирования и редактирования геометрии. Для вывода на растровый экран векторное описание в конечном счёте \textbf{растеризуется}.

\subsection{Геометрические преобразования}

Основные преобразования --- перенос, масштабирование, вращение и проецирование. В трёхмерной графике удобно использовать \textbf{однородные координаты}
\[
p=(x,y,z,1)^T,
\]
тогда аффинное преобразование записывается как
\[
p'=Mp.
\]
Например,
\[
T=
\begin{pmatrix}
1&0&0&t_x\\
0&1&0&t_y\\
0&0&1&t_z\\
0&0&0&1
\end{pmatrix},
\qquad
R_z=
\begin{pmatrix}
\cos\varphi&-\sin\varphi&0&0\\
\sin\varphi& \cos\varphi&0&0\\
0&0&1&0\\
0&0&0&1
\end{pmatrix}.
\]
Преобразования композиционируются умножением матриц; порядок существенен.

Типичная цепочка систем координат:
\[
\text{локальные}
\rightarrow
\text{мировые}
\rightarrow
\text{координаты камеры}
\rightarrow
\text{clip}
\rightarrow
\text{NDC}
\rightarrow
\text{экран}.
\]
Для вершины обычно пишут
\[
p_{\mathrm{clip}}=PVMp_{\mathrm{local}},
\]
где $M$ --- матрица модели, $V$ --- вида, $P$ --- проекции. При перспективной проекции после преобразования выполняется перспективное деление
\[
x_{\mathrm{NDC}}=\frac{x_c}{w_c},\quad
 y_{\mathrm{NDC}}=\frac{y_c}{w_c},\quad
 z_{\mathrm{NDC}}=\frac{z_c}{w_c}.
\]

\subsection{Графический конвейер}

\textbf{Графический конвейер} --- последовательность этапов преобразования геометрического описания сцены в кадр. Упрощённо:
\[
\begin{split}
&\text{вершины}
\rightarrow \text{вершинная обработка}
\rightarrow \text{сборка примитивов}
\rightarrow \text{отсечение}\\
&\rightarrow \text{растеризация}
\rightarrow \text{обработка фрагментов}
\rightarrow \text{кадр}.
\end{split}
\]
Основным примитивом 3D-растеризации обычно является треугольник. Вершина может содержать координаты, нормаль, цвет, текстурные координаты и другие атрибуты.

\textbf{Растеризация} преобразует геометрический примитив в набор фрагментов, соответствующих потенциальным пикселям. Для определения видимой поверхности обычно применяется \textbf{$Z$-буфер}: для каждого пикселя хранится глубина ближайшего фрагмента.

Цвет фрагмента определяется материалом, освещением и текстурами. Например, простая диффузная модель Ламберта содержит множитель
\[
\max(0,\mathbf n\cdot\mathbf l),
\]
где $\mathbf n$ --- нормаль, $\mathbf l$ --- направление на источник света.

Растеризация --- не единственный способ рендеринга. При \textbf{трассировке лучей} из камеры строятся лучи
\[
r(t)=o+td
\]
и ищутся их пересечения со сценой; этот подход естественно описывает тени, отражения и преломления, но обычно вычислительно дороже.

\subsection{Графические библиотеки и API}

\textbf{Графический API} предоставляет программе стандартный интерфейс к графическому оборудованию. Примеры: OpenGL, Vulkan, Direct3D, Metal, WebGL. Через API создаются и заполняются буферы, загружаются текстуры, задаются состояния конвейера и формируются команды рисования.

Современный конвейер программируем. Его стадии задаются \textbf{шейдерами}:
\begin{itemize}
\item вершинный шейдер преобразует вершины, например $p_{\mathrm{clip}}=PVMp$;
\item фрагментный шейдер вычисляет цвет фрагмента с учётом материалов, освещения и текстур;
\item в современных API могут использоваться также геометрические, тесселяционные и вычислительные шейдеры.
\end{itemize}

Высокоуровневые библиотеки и системы визуализации скрывают значительную часть низкоуровневой работы с GPU. В пакетах трёхмерной графики типичный цикл включает моделирование геометрии, материалы и текстуры, освещение и рендеринг; растеризационные движки ориентированы на интерактивность, а трассировочные --- на более физически точное моделирование света.

\subsection{Научная визуализация}

В вычислительном эксперименте графика используется для анализа полей и данных. Для скалярных полей применяют графики, карты цветов, линии уровня
\[
f(x,y)=C,
\]
сечения и изоповерхности
\[
f(x,y,z)=C.
\]
Для векторных полей используют стрелочные диаграммы, линии тока и траектории. В механике визуализируют сетки, поля температур, давления, скоростей, напряжений и деформаций.

\subsection{Что важно сказать на экзамене}

Минимальная логика ответа:
\[
\boxed{
\text{растровое/векторное представление}
\rightarrow
\text{матричные преобразования}
\rightarrow
\text{графический конвейер}
\rightarrow
\text{API и шейдеры}
}.
\]
Нужно уметь объяснить роль однородных координат, перспективного преобразования, растеризации и $Z$-буфера и назвать хотя бы одну графическую библиотеку/API.

\newpage
